\section{Related Work}
% \subsection{LLM에서의 다국어 추론}
% 거대언어모델(LLM)은 Chain-of-Thought (CoT) 프롬프팅~\citep{cot-neurips2022}을 통해 강력한 추론 능력을 보여주지만, 이러한 능력은 언어에 따라 상당한 차이를 보이며 비영어권 언어에서는 지속적으로 저조한 성능을 보인다~\citep{shi2022language,mega-emnlp2023}.

% 이러한 격차를 해소하기 위해, 최근 연구들은 다국어 데이터에 대한 지도 미세 조정(SFT)을 활용해 왔다. \citet{lai2024mcot}은 다국어 수학 추론 데이터셋인 mCoT-MATH를 소개했으며, \citet{chai2025xcot}는 교차 언어 지시 튜닝(cross-lingual instruction tuning)을 제안했다. 그러나 SFT 접근법은 저자원 언어에 대한 데이터 희소성 문제로 인해 한계가 있다~\citep{costa2022no}.

% 프롬프트 기반 방법론은 추가적인 훈련 없이 LLM의 잠재된 지식을 활용함으로써 대안을 제시한다~\citep{brown2020language}. 퓨샷 인컨텍스트 러닝(Few-shot in-context learning)~\citep{lin2022few}과 다양한 CoT 변형 기법들~\citep{kojima2022large,zhouleast}은 여러 언어에 걸쳐 유연성을 제공하지만, 영어 중심의 성능 격차는 여전히 존재한다.

% \subsection{번역 기반 추론 대 대상 언어 추론}
% 다국어 추론은 두 가지 패러다임 사이의 근본적인 상충 관계(trade-off)에 직면해 있다. 번역 기반 접근법은 추론 전에 입력값을 영어로 변환함으로써 LLM의 더 강력한 영어 능력을 활용한다. Cross-Lingual Prompting (CLP)~\citep{clp-emnlp2023}과 Cross-Lingual Thought (XLT)~\citep{huang2023not}는 영어를 매개로 한(English-pivoted) 추론을 통해 성능 향상을 입증했다.

% 그러나 번역 기반 방법론은 환각(hallucinations), 의미적 손실(semantic loss)~\citep{dale2023halomi}, 계산 오버헤드 증가, 그리고 오류 전파(error propagation)~\citep{zhu2024question}와 같은 내재적인 단점들을 안고 있다. 게다가 영어로만 구성된 출력은 비영어권 사용자에게 해석 가능성(interpretability)을 떨어뜨린다~\citep{multilingual-prompting-emnlp2025}.

% 반면, 대상 언어(target-language)로 추론을 수행하는 것은 의미적 충실도를 보존하고 사용자가 추론 과정을 검증할 수 있게 하지만~\citep{shi2022language,reason-lang-emnlp2025}, LLM의 상대적으로 약한 비영어 능력으로 인해 종종 낮은 정확도를 초래한다~\citep{huang2023not}. 

% \subsection{구조적 추론 접근법}
% 최근 추론 전략의 발전은 단순한 선형적 생성을 넘어 구조적 분해(structured decomposition)로 나아가고 있다. Plan-and-Solve~\citep{wang2023plan}와 Program-of-Thought~\citep{chen2022program}는 문제를 계획이나 코드로 분해함으로써 신뢰성을 향상시킨다. 이와 유사하게, Skeleton-of-Thought~\citep{ning2023skeleton}는 답변의 간결한 개요(outline)를 구성하는데, 이는 본래 병렬 디코딩을 통한 추론 가속화를 목표로 제안되었다.

% 우리는 다국어 환경에서 이러한 스켈레톤 구조의 효용성을 재해석한다. 복잡한 그래프 기반 표현에 의존하는 Structured-of-Thought~\citep{qi2025sot}와 달리, 우리의 접근 방식은 가벼운(lightweight) 영어 텍스트 스켈레톤을 추론 가이드로 활용한다. 이 스켈레톤은 모델의 강력한 영어 추론 논리를 대상 언어(target language)로 전이하는 중간 발판(scaffold) 역할을 수행한다. (영어에서의) 논리적 계획과 (대상 언어에서의) 언어적 서술(articulation)을 분리함으로써, 우리의 방식은 전체 텍스트 번역에 따르는 계산 비용 없이 높은 정확도와 원어민 수준의 해석 가능성을 달성한다.

\begin{comment}
\subsection{English-Pivot Reasoning} 대규모 언어 모델(LLM)은 Chain-of-Thought (CoT) 프롬프팅~\citep{cot-neurips2022}을 통해 강력한 추론 능력을 보여주지만, 이러한 능력은 언어에 따라 상당한 차이를 보이며 비영어권 언어에서는 지속적으로 저조한 성능을 나타낸다~\citep{shi2022language,mega-emnlp2023}.  특히, 최근 연구들은 다국어 데이터에 대해 supervised fine-tuning(SFT)방식을 도입하였지만, 저자원 언어에 대한 데이터 희소성 문제로 인해 한계가 있다~\citep{costa2022no,lai2024mcot,chai2025xcot}. 이러한 성능 격차를 완화하기 위해 영어를 피벗(pivot) 언어로 활용하는 프롬프팅 전략들이 제안되었다. Cross-Lingual Prompting (CLP)~\citep{clp-emnlp2023}과 Cross-Lingual Thought (XLT)~\citep{xlt-emnlp2023}는 비영어 질문을 영어로 변환하여 의미를 파악하게 하거나 영어로 추론을 수행함으로써 성능 향상을 입증하였다. 또한, \citet{trans-allneed-2025naacl}은 영어 중심(English-centric) LLM에서 비영어 질문을 영어로 번역하는 것이 다양한 이해 및 추론 작업에 효과적임을 보여주었다. 이는 Multilingual LLM(MLLM)이 비영어 입력에 대해 겪는 질문 이해의 어려움을 영어의 우수한 언어 능력을 통해 보완할 수 있음을 시사한다. 나아가 구조화된 프롬프팅을 통한 개선 시도도 이루어졌다. \citet{qi2025sot}는 영어를 활용하여 문제 해결에 필요한 변수를 추출하고 각 언어별 지식을 결합하는 Structured-of-Thought를 제안하였으며, \citet{li2024eliciting}는 코드 형태의 구조로 추론 과정을 유도하여 모델이 단계적으로 문제를 해결하도록 설계하였다. 이러한 접근법들은 공통적으로 영어의 강력한 추론 능력을 활용하여 multilingual reasoning gap을 줄이는 데 주력하였다.

\subsection{The Challenge of Target-Language Reasoning} 선행 연구들은 비영어권 사용자의 관점, 특히 결과의 해석 가능성(interpretability)과 언어적 특성을 충분히 고려하지 못했다는 한계가 있다. 물론 Target Language로 추론을 수행하는 것은 사용자가 과정을 직접 검증할 수 있게 하며, 문화적 맥락이 중요한 과제에서는 영어보다 더 적합한 결과를 보이기도 한다~\citep{trans-allneed-2025naacl,multilingual-prompting-emnlp2025}. 그러나 수학이나 복잡한 논리 구조가 요구되는 과제에서 target language CoT는 여전히 영어에 비해 저조한 성능을 보인다. \citet{shi2022language}와 \citet{wang2025polymath}는 비영어 언어로 추론 시 성능 하락이 발생함을 보고하였으며, 특히 \citet{long-cot-2025arxiv}는 이러한 하락이 개념 적용(conceptual application) 단계에서의 오류에 기인함을 실증적으로 밝혔다. 이는 MLLM이 비영어 환경에서 도메인 지식과 논리적 구조를 온전히 활용하는 데 여전히 어려움이 있음을 시사한다.

이에 본 논문에서는 문제 구조화와 핵심 개념으로 구성된 Skeleton을 활용한 Skeleton-Guided Chain-of-Thought(CoT)를 제안한다. 영어로 생성된 Skeleton을 논리적 가이드로 사용함으로써, target language CoT에서 발생하는 개념적 오류와 논리적 불안정성을 완화하고, 동시에 비영어권 사용자에게 필요한 해석 가능성을 유지한다.
\end{comment}


\subsection{English-Pivot Reasoning}
대규모 언어 모델은 Chain-of-Thought (CoT) 프롬프팅~\citep{cot-neurips2022}을 통해 강력한 추론 능력을 보여주었으나, 이러한 능력은 언어에 따라 상당한 편차를 보인다. 다수의 선행 연구들은 비영어권 언어에서 추론 성능이 지속적으로 저하됨을 보고하였으며, 이러한 격차는 특히 수학적·논리적 추론 과제에서 두드러지게 나타난다~\citep{shi2022language, mega-emnlp2023, wang2025polymath}.

이러한 성능 차이를 완화하기 위해 영어를 피벗(pivot) 언어로 활용하는 프롬프팅 전략이 제안되었다. CLP~\citep{clp-emnlp2023}와 XLT~\citep{xlt-emnlp2023}는 비영어 질문을 영어로 변환하여 의미를 파악하게 하거나 추론 과정을 영어로 수행함으로써 성능 향상을 입증하였다. 후속 연구들 또한 비영어 입력을 영어로 번역하여 처리하는 것이 다양한 이해 및 추론 작업에 효과적임을 보여주었다~\citep{trans-allneed-2025naacl}.

그러나 이러한 접근법은 최종 답변이 영어로 생성될 경우 비영어권 사용자가 결과를 직관적으로 이해하거나 검증하기 어려워, 실용적인 측면에서 interpretability이 떨어진다는 한계를 가진다~\citep{multilingual-prompting-emnlp2025}.

\subsection{Structured Reasoning and Language Configuration} 
단순한 언어 변환을 넘어, 추론 과정을 명시적으로 구조화하여 복잡한 문제를 해결하려는 시도 또한 활발히 이루어졌다. Skeleton-of-Thought~\citep{ning2023skeleton}나 Plan-and-Solve~\citep{wang2023plan}와 같은 연구는 문제 해결을 위한 뼈대(skeleton)나 계획을 먼저 생성하는 것이 모델의 논리적 일관성을 높이고 오류를 줄이는 데 효과적임을 보였다. 다국어 환경에서도 \citet{qi2025sot}는 영어로 변수와 지식을 구조화하는 방식을 제안하였고, \citet{li2024eliciting}는 코드(Code) 형식을 빌려 단계적 추론을 유도하였다.

하지만 기존의 skeleton과 같은 구조화된 방법론들은 정보를 모두 영어로 제공하거나, 입력부터 출력까지의 언어 구성을 단일하게 고정하여 실험하였다는 한계가 있다. 이를 완화하기 위한 수단으로서 스켈레톤 언어와 최종 답변 언어의 조합이 미치는 영향을 체계적으로 분석하지 않았다.

선행 연구들은 Target Language로 직접 추론(CoT)을 수행할 때 발생하는 개념적 적용 오류(conceptual application error)~\citep{long-cot-2025arxiv}를 지적하면서도, 

본 연구는 이러한 공백을 다루며, 영어 스켈레톤의 논리적 강점을 활용하면서도 Target Language 답변의 효용성을 유지하기 위한 최적의 언어 구성을 탐구한다는 점에서 기존 연구들과 차별화된다.

% 그러나 최종 답변이 영어로 생성될 경우 비영어권 사용자가 결과를 직관적으로 이해하고 검증하기 어렵다는 실용적 한계(interpretability challenge)를 갖는다. 더욱이 기존 연구들은 대부분 입력 언어 또는 추론 언어와 같은 \textit{단일 언어 요소}의 변환에 초점을 맞추어 왔으며, 입력 질의 언어, 추론 구조를 제공하는 중간 표현(스켈레톤), 그리고 실제 추론이 수행되는 언어 간의 \textbf{조합적 상호작용}을 독립적인 설계 변수로 분석한 연구는 제한적이다. 특히 저자원 언어 환경에서 이러한 언어 조합이 추론 성능에 미치는 영향에 대해서는 체계적인 분석이 이루어지지 않았다.